\section{Estimación del Alcance}

El presente proyecto tiene como objetivo ofrecer un servicio de detección automática de lenguaje ofensivo en redes sociales y plataformas digitales, haciendo uso de modelos basados en arquitecturas Transformers. El servicio está diseñado para ser integrado principalmente por medios digitales ---como portales de noticias y diarios en línea--- que actualmente han restringido o eliminado la funcionalionalidad de comentarios por dificultades en su moderación. Asimismo, se proyecta como una solución viable para medios televisivos que requieren mostrar comentarios del público en transmisiones en vivo, garantizando previamente su moderación.

El valor agregado del servicio radica en su capacidad para filtrar, en tiempo real, contenido ofensivo o inapropiado proveniente de usuarios, asegurando que las interacciones públicas cumplan con estándares de respeto y convivencia digital. Esta funcionalidad permitirá a los medios:

\begin{itemize}
    \item Reactivar espacios de comentarios sin temor al contenido ofensivo.
    
    \item Mostrar en vivo interacciones del público en emisiones televisivas con moderación automática.
    
    \item Ofrecer entornos seguros en sus propias plataformas para fomentar la participación del público.
\end{itemize}

A futuro, se contempla la expansión del servicio hacia canales de transmisión en línea (como Kick) y comunidades virtuales (como foros especializados o aplicaciones de mensajería grupal), donde el volumen y la inmediatez de los mensajes hacen indispensable una moderación automatizada.

\section{Costos de Operación}

\subsection{Mano de Obra Directa}

Para la etapa inicial del proyecto se contempla una estructura mínima de personal. Esta decisión busca optimizar recursos sin comprometer la calidad funcional del MVP. Conforme se validen hipótesis del mercado y aumente la base de usuarios, se contempla incorporar personal adicional.

Partiendo desde un equipo de desarrollo conformado por:

\begin{itemize}
    \item \textbf{Desarrollador Full Stack:} Se encarga tanto del backend (API de detección) como del frontend (extensión, panel), integrando la funcionalidad del modelo y las interfaces.
\end{itemize}
